
Compilation of \openshmem GPU programs requires both the \openshmem
implementation's compile wrapper and the GPU compiler toolchain. The exact
commands depend on the implementation and GPU vendor.

\subsection*{Compilation}

A typical compilation workflow proceeds as follows:

\begin{enumerate}
\item GPU device code is compiled using the vendor compiler (e.g., \texttt{nvcc},
  \texttt{hipcc}, \texttt{icpx -fsycl}).
\item The resulting objects are linked with the \openshmem library that supports
  the GPU extensions.
\item Both \texttt{shmem.h} and \texttt{shmemg.h} are made available in the
  include path.
\end{enumerate}

Example (CUDA):
\begin{lstlisting}[language=bash,style=SourceListingsDefault]
nvcc -c gpu_kernel.cu -o gpu_kernel.o
oshcc -c host_code.c -o host_code.o
oshcc gpu_kernel.o host_code.o -o my_program -lcudart
\end{lstlisting}

\subsection*{Execution}

Programs are launched using the \openshmem launcher (e.g., \texttt{oshrun}):
\begin{lstlisting}[language=bash,style=SourceListingsDefault]
export SHMEM_DEVICE_SYMMETRIC_SIZE=256m
oshrun -n 4 ./my_program
\end{lstlisting}

The \ENVVAR{SHMEM\_DEVICE\_SYMMETRIC\_SIZE} environment variable controls the
size of the \ac{GPU} symmetric heap allocated per \ac{PE}.

\clearpage

